% DERIVED from formal_layer.json — read-only, do not edit.
\begin{definitionv}{synth_2}[Cell-specific treatment propensity]
For a probability law \(P\) of \((X,A,Y)\) on \(\{1,\ldots,d\}\times\{0,1\}^2\), let \(p_j(P)=P(X=j)\). Define the treatment propensity in cell \(j\) by
\[
e_j(P)=
\begin{cases}
\Pr_P(A=1\mid X=j)
=\dfrac{P(X=j,A=1)}{p_j(P)},&p_j(P)>0,\\[6pt]
\dfrac12,&p_j(P)=0.
\end{cases}
\]
\end{definitionv}

\begin{assumptionv}{ass:overlap}[Public overlap restriction]
For every \(j\) with \(p_j(P)>0\),
\[
e_j(P)\in[\epsilon,1-\epsilon].
\]
\end{assumptionv}

\begin{definitionv}{synth_1}[Covariate and treatment marginal]
The treatment–covariate marginal \(P_{XA}\) is the probability law obtained from \(P\), the probability law of complete records \((X,A,Y)\), by the projection \((X,A,Y)\mapsto(X,A)\). Here \(X\) takes values in a finite alphabet of size \(d\), for a nonnegative integer \(d\), and \(A,Y\in\{0,1\}\). Its cell-probability table is
\[
P_{XA}(j,a)
=
\sum_{y\in\{0,1\}} P(X=j,A=a,Y=y),
\]
for each covariate value \(j\) and treatment value \(a\in\{0,1\}\).
\end{definitionv}

\begin{assumptionv}{ass:data-product}[Two-channel product sampling]
The observed data \(\mathcal D\) have distribution
\[
\mathcal D\sim P^{\otimes n}\otimes P_{XA}^{\otimes m}.
\]
\end{assumptionv}

\begin{definitionv}{synth_5}[Probability laws]
For a measurable space \(\mathcal Z\), let \(\mathsf{Prob}(\mathcal Z)\) denote the class of probability measures on \(\mathcal Z\).
\end{definitionv}

\begin{definitionv}{def:model}[Observable law class \(\mathcal M_{d,\epsilon}\)]
The observable law class is
\[
\mathcal M_{d,\epsilon}
=
\left\{
P\in\mathsf{Prob}\bigl(\{1,\ldots,d\}\times\{0,1\}^2\bigr):
P\text{ satisfies }\cref{obj:ass:overlap}
\right\},
\qquad d\ge2,\quad 0<\epsilon\le\tfrac14.
\]
This is model (3) of Zeng et al., restricted to the stated range of the public overlap floor \(\epsilon\).
\end{definitionv}

\begin{definitionv}{def:target}[Outcome-contrast estimand \(\tau(P)\)]
The estimand is
\[
\tau(P)
=
\sum_{j=1}^{d}p_j(P)\{\mu_{1j}(P)-\mu_{0j}(P)\},
\]
with summands for null covariate cells defined to equal zero.
\end{definitionv}

\begin{definitionv}{synth_3}[Independent uniform randomization]
Let \(\mathrm{Uniform}[0,1]\) denote the probability law on \([0,1]\) assigning each Borel set its Lebesgue measure. The estimator randomization seed \(U\) has this law and is independent of the observed data \(\mathcal D=(\mathcal L,\mathcal V)\), where \(\mathcal L\) and \(\mathcal V\) are the ordered complete-record and auxiliary samples, respectively.
\end{definitionv}

\begin{definitionv}{def:risk}[Minimax risk $R(n,m,d,\epsilon)$]
The minimax squared-error risk is
\[
R(n,m,d,\epsilon)
=
\inf_T\sup_{P\in\mathcal M_{d,\epsilon}}
\mathbb E_{P^{\otimes n}\otimes P_{XA}^{\otimes m}\otimes\mathrm{Uniform}[0,1]}
\bigl[(T(\mathcal D,U)-\tau(P))^2\bigr].
\]
The infimum ranges over all Borel estimation rules on the original observed records, including rules randomized through the independent seed \(U\). Clipping under squared-error loss permits restriction to rules taking values in \([-1,1]\).
\end{definitionv}

\begin{assumptionv}{ass:consistency}[Potential-outcome consistency condition]
Under the potential-outcome extension \(H\),
\[
Y=AY_1+(1-A)Y_0
\]
holds \(H\)-almost surely.
\end{assumptionv}

\begin{assumptionv}{ass:exchangeability}[Conditional treatment exchangeability]
Under the potential-outcome extension \(H\),
\[
(Y_0,Y_1)\perp\!\!\!\perp A\mid X.
\]
\end{assumptionv}

\begin{definitionv}{def:known-risk}[Exact-marginal risk $R^{\mathrm K}(n,d,\epsilon)$]
Define the minimax squared-error risk with the exact treatment--covariate marginal supplied by
\[
R^{\mathrm K}(n,d,\epsilon)
=
\inf_{T^{\mathrm K}}\sup_{P\in\mathcal M_{d,\epsilon}}
\mathbb E_{P^{\otimes n}\otimes\mathrm{Uniform}[0,1]}
\left[
\bigl(T^{\mathrm K}(\mathcal L,P_{XA},U)-\tau(P)\bigr)^2
\right].
\]
Here \(\mathcal L\) is the complete-record component of \(\mathcal D\), and the infimum ranges over Borel estimation rules receiving \(\mathcal L\), the exact marginal table \(P_{XA}\), and the independent seed \(U\).
\end{definitionv}

\begin{definitionv}{def:rate-handle}[Numerical rate $r(n,m,d,\epsilon)$]
Define
\[
r(n,m,d,\epsilon)
=
\min\left\{
1,\frac1S+\left(\frac d{N\epsilon\ell}\right)^2
\right\},
\qquad
N=n+m,\quad S=n\epsilon,\quad
\ell=\log(\exp(1)+S).
\]
All denominators are strictly positive on the public index range, and \(r(n,m,d,\epsilon)\) is a positive numerical function.
\end{definitionv}

\begin{definitionv}{synth_4}[Clipping to the target range]
For every \(x\in\mathbb R\), define
\[
\operatorname{clip}_{[-1,1]}(x)=\max\{-1,\min\{x,1\}\}.
\]
\end{definitionv}

\begin{algorithmv}{def:estimator-handle}[Prefix-average estimator $\widehat\tau$]
The estimator takes the ordered complete and auxiliary records \(\mathcal D\), the independent seed \(U\), and the public parameters \(n,m,d,\epsilon\) as inputs.
\begin{enumerate}
\item Set
\[
S=n\epsilon.
\]
If \(S<\exp(4096)\), output
\[
\widehat\tau(\mathcal D,U)=0.
\]
For \(S\ge\exp(4096)\), proceed through the remaining steps.

\item Define the pool lengths and tuning parameters by
\[
\begin{aligned}
h_0&=\lfloor n/3\rfloor,
&
h_p&=\lfloor(n-h_0)/2\rfloor,
&
h_f&=n-h_0-h_p,\\
M_p&=h_p+\lfloor m/2\rfloor,
&
M_f&=h_f+m-\lfloor m/2\rfloor,\\
u&=h_0/8,
&
t_p&=M_p/8,
&
t&=M_f/8,\\
L&=\lfloor\log S/1024\rfloor,
&
B&=\frac{2^{20}L}{\min\{t_p,t\}},
&
k_0&=\lfloor t_pB/4\rfloor.
\end{aligned}
\]
Use the first \(h_0\) complete records as the outcome pool. Form the pilot pool by concatenating the treatment--covariate projections of the next \(h_p\) complete records with the first \(\lfloor m/2\rfloor\) auxiliary records. Form the factorial pool by concatenating the remaining complete-record projections with the remaining auxiliary records. Preserve the stated order within each pool.

\item Define the Chebyshev polynomials and the reciprocal-approximation coefficients by
\[
\begin{aligned}
T_0(x)&=1,\qquad T_1(x)=x,\\
T_{h+1}(x)&=2xT_h(x)-T_{h-1}(x),\qquad h\ge1,\\
E_L(z)&=\frac{1-T_L(1-2z)}{2L^2z},\\
G_L(z)&=\frac{1-E_L(z)}z
=\sum_{h=0}^{L-2}g_hz^h.
\end{aligned}
\]
Use the removable extensions at zero in the definitions of \(E_L\) and \(G_L\).

\item For each prefix triple satisfying
\[
0\le h'\le h_0,\qquad
0\le h''\le M_p,\qquad
0\le h'''\le M_f,
\]
write the ordered outcome-pool observations as
\((X_i^{\mathrm o},A_i^{\mathrm o},Y_i^{\mathrm o})\), the pilot-pool observations as
\((X_i^{\mathrm p},A_i^{\mathrm p})\), and the factorial-pool observations as
\((X_i^{\mathrm f},A_i^{\mathrm f})\). For \(a\in\{0,1\}\) and \(j\in\{1,\ldots,d\}\), define
\[
\begin{aligned}
Z_{aj}
&=\sum_{i=1}^{h'}
\mathbf1\{X_i^{\mathrm o}=j,\ A_i^{\mathrm o}=a,\ Y_i^{\mathrm o}=1\},\\
J_{aj}
&=\sum_{i=1}^{h''}
\mathbf1\{X_i^{\mathrm p}=j,\ A_i^{\mathrm p}=a\},\\
K_{aj}
&=\sum_{i=1}^{h'''}
\mathbf1\{X_i^{\mathrm f}=j,\ A_i^{\mathrm f}=a\}.
\end{aligned}
\]
Define the falling factorials of an integer count \(x\) by
\[
(x)_0=1,\qquad
(x)_h=\prod_{r=0}^{h-1}(x-r),\qquad h\ge1.
\]

\item For each prefix triple, define
\[
\begin{aligned}
\widehat G_{aj}
&=\sum_{h=0}^{L-2}g_h\frac{(K_{aj})_h}{(tB)^h},\\
P_{aj}^{\mathrm{pol}}
&=\frac{Z_{aj}}u
\left(1+\frac{K_{1-a,j}\widehat G_{aj}}{tB}\right),\\
H_{aj}
&=\frac{Z_{aj}}u
\left(1+\frac{K_{1-a,j}}{K_{aj}+1}\right),\\
V_{aj}
&=\mathbf1\{J_{aj}\le k_0\}P_{aj}^{\mathrm{pol}}
+\mathbf1\{J_{aj}>k_0\}H_{aj},\\
F(h',h'',h''';\mathcal D)
&=\operatorname{clip}_{[-1,1]}
\left(\sum_{j=1}^d(V_{1j}-V_{0j})\right).
\end{aligned}
\]

\item Output the deterministic finite average
\[
\widehat\tau(\mathcal D,U)
=
\sum_{h'=0}^{h_0}\sum_{h''=0}^{M_p}\sum_{h'''=0}^{M_f}
\exp(-u-t_p-t)
\frac{u^{h'}}{h'!}
\frac{t_p^{h''}}{h''!}
\frac{t^{h'''}}{h'''!}
F(h',h'',h''';\mathcal D).
\]
Poisson prefix combinations exceeding a pool length have output zero. The estimator is independent of \(U\). Every denominator is a strictly positive public quantity or an integer count plus one, including for empty observed cells; all quantities in the construction are determined by the observed data and public parameters.
\end{enumerate}
\end{algorithmv}

\begin{theoremv}{thm:uniform-annotation-frontier}[Uniform annotation frontier]
There exist numerical constants \(0<c\le C<\infty\), independent of all public indices, such that the following holds for every experiment with:
\begin{itemize}
\item \textbf{(Sample sizes.)} Integers \(n\ge1\) and \(m\ge0\).
\item \textbf{(Covariate alphabet.)} An integer \(d\ge2\).
\item \textbf{(Overlap floor.)} A real parameter \(0<\epsilon\le1/4\).
\end{itemize}
The deterministic estimator \(\widehat\tau(\mathcal D,U)\) constructed in \cref{obj:def:estimator-handle} is Borel measurable, ignores the randomization seed \(U\), and takes values in \([-1,1]\) for every data realization. With \(R(n,m,d,\epsilon)\) the minimax squared-error risk from \cref{obj:def:risk} and \(r(n,m,d,\epsilon)\) the rate from \cref{obj:def:rate-handle},
\[
c\,r(n,m,d,\epsilon)
\le R(n,m,d,\epsilon)
\le \sup_{P\in\mathcal M_{d,\epsilon}}
\mathbb E_P\!\left[
\bigl(\widehat\tau(\mathcal D,U)-\tau(P)\bigr)^2
\right]
\le C\,r(n,m,d,\epsilon).
\]
Here \(\mathcal M_{d,\epsilon}\) is the observable model class in \cref{obj:def:model}, \(\tau(P)\) is the observable target in \cref{obj:def:target}, and the expectation uses the experiment specified in \cref{obj:def:risk}.

Moreover, for every integer \(d\ge2\), every \(0<\epsilon\le1/4\), and every observable law \(P\in\mathcal M_{d,\epsilon}\), there exists a joint law \(H\) of the observed coordinates and binary potential outcomes \(Y_0,Y_1\) whose observed marginal is \(P\) and which satisfies \cref{obj:ass:consistency,obj:ass:exchangeability}. Every such law \(H\) satisfies
\[
\mathbb E_H[Y_1-Y_0]=\tau(P).
\]
\end{theoremv}

\begin{theoremv}{thm:rare-label-floor}[Uniform rare-label risk floors]
There is a numerical constant \(c>0\), independent of all public indices, such that, for every \(n,m,d\in\mathbb N\) satisfying
\begin{itemize}
\item \textbf{(Complete-record sample size.)} \(n\ge 1\).
\item \textbf{(Covariate alphabet.)} \(d\ge 2\).
\item \textbf{(Overlap parameter.)} \(0<\epsilon\le 1/4\).
\end{itemize}
the minimax risks in \cref{obj:def:risk,obj:def:known-risk} satisfy
\[
R(n,m,d,\epsilon)\ge c\,b(n,\epsilon),
\qquad
R^{\mathrm K}(n,d,\epsilon)\ge c\,b(n,\epsilon),
\qquad
b(n,\epsilon)=\min\{1,(n\epsilon)^{-1}\}.
\]
The first bound applies to the infimum over all clipped Borel randomized original-record rules specified in \cref{obj:def:risk}.

Moreover, for every \(d\in\mathbb N\), overlap parameter \(\epsilon\in\mathbb R\), and separation \(\delta\in\mathbb R\) satisfying
\begin{itemize}
\item \textbf{(Covariate alphabet.)} \(d\ge 2\).
\item \textbf{(Overlap parameter.)} \(0<\epsilon\le 1/4\).
\item \textbf{(Separation.)} \(0<\delta\le 1/8\).
\end{itemize}
the two complete-record laws \(P_+\) and \(P_-\) specified by
\[
\begin{aligned}
p_j(P_\pm)&=\frac12,
&
e_j(P_\pm)&=\epsilon,
&
\mu_{0j}(P_\pm)&=\frac12,
&
\mu_{1j}(P_\pm)&=\frac12\pm\delta,
&& j=1,2,\\
p_j(P_\pm)&=0,
&&&&&&&& j=3,\ldots,d
\end{aligned}
\]
belong to \(\mathcal M_{d,\epsilon}\) of \cref{obj:def:model}. Their targets, defined in \cref{obj:def:target}, and treatment–covariate marginals satisfy
\[
\tau(P_+)=\delta,
\qquad
\tau(P_-)=-\delta,
\qquad
(P_+)_{XA}=(P_-)_{XA}.
\]
At \(\delta=1/8\), their one-record Kullback--Leibler divergence is exactly
\[
D(P_+\Vert P_-)
=
\sum_{z:\,P_+(z)>0}
P_+(z)\log\frac{P_+(z)}{P_-(z)}
=
\frac{\epsilon}{4}\log\frac53,
\]
where \(z\) ranges over complete-record atoms and \(\log\) is the natural logarithm.
\end{theoremv}

\begin{propositionv}{prop:benchmark-recovery}[Benchmark recovery and oracle comparison]
Let \(r(n,m,d,\epsilon)\), \(R(n,m,d,\epsilon)\), and \(R^{\mathrm K}(n,d,\epsilon)\) be as in \cref{obj:def:rate-handle,obj:def:risk,obj:def:known-risk}. Define the rare-label and fixed-overlap benchmark functions by
\[
b(n,\epsilon)=\min\left\{1,\frac{1}{n\epsilon}\right\},
\qquad
f(n,m,d)=\min\left\{1,\frac{1}{n}
+\frac{d^2}{(n+m)^2[\log(\exp(1)n)]^2}\right\}.
\]
The following statements hold.
\begin{enumerate}
\item For every fixed \(0<\epsilon\le 1/4\), there exist constants \(a,B_c>0\), depending only on \(\epsilon\), such that for all integers \(n\ge1\), \(m\ge0\), and \(d\ge2\),
\[
a f(n,m,d)\le r(n,m,d,\epsilon)\le B_c f(n,m,d).
\]

\item There exist numerical constants \(c,C>0\) such that for all integers \(n\ge1\), \(m\ge0\), and \(d\ge2\), and every \(0<\epsilon\le1/4\),
\[
c b(n,\epsilon)\le R(n,m,2,\epsilon)\le C b(n,\epsilon),
\qquad
c b(n,\epsilon)\le R^{\mathrm K}(n,d,\epsilon)\le C b(n,\epsilon).
\]

\item For all integers \(n\ge1\), \(m\ge1\), and \(d\ge2\), and every \(0<\epsilon\le1/4\),
\[
0\le R(n,m,d,\epsilon)-R^{\mathrm K}(n,d,\epsilon)
\le \min\left\{1,6(n+2)\sqrt{\frac{2d}{m}}\right\}.
\]

\item For every fixed integer \(n\ge1\), integer \(d\ge2\), and \(0<\epsilon\le1/4\), the limit along integer auxiliary sample sizes satisfies
\[
\lim_{m\to\infty}R(n,m,d,\epsilon)=R^{\mathrm K}(n,d,\epsilon).
\]

\item For every real \(M_0>0\), there exists \(\kappa>0\), depending only on \(M_0\), such that for all integers \(n\ge1\), \(m\ge0\), and \(d\ge2\), and every \(0<\epsilon\le1/4\) satisfying \(n\epsilon\le M_0\),
\[
R(n,m,d,\epsilon)\ge\kappa.
\]
\end{enumerate}
\end{propositionv}

\begin{algorithmv}{def:baseline}[Inverse-count baseline $T^{\mathrm B}$]
The baseline takes the ordered complete and auxiliary records \(\mathcal D\) and the public parameters \(n,m,d,\epsilon\) as inputs.
\begin{enumerate}
\item If \(n<6\), output
\[
T^{\mathrm B}(\mathcal D)=0.
\]
For \(n\ge6\), proceed through the remaining steps.

\item Define the pool lengths and Poisson prefix means by
\[
h_{\mathrm o}=\lfloor n/2\rfloor,\qquad
h_{\mathrm m}=n-h_{\mathrm o}+m,\qquad
u_{\mathrm B}=h_{\mathrm o}/8,\qquad
t_{\mathrm B}=h_{\mathrm m}/8.
\]
Use the first \(h_{\mathrm o}\) complete records as the outcome pool. Form the marginal pool by concatenating the treatment--covariate projections of the remaining complete records with all auxiliary records, preserving their order.

\item Write the outcome-pool observations as
\((X_i^{\mathrm o},A_i^{\mathrm o},Y_i^{\mathrm o})\) and the marginal-pool observations as
\((X_i^{\mathrm m},A_i^{\mathrm m})\). For each prefix pair satisfying
\[
0\le h'\le h_{\mathrm o},\qquad
0\le h''\le h_{\mathrm m},
\]
and for \(a\in\{0,1\}\), \(j\in\{1,\ldots,d\}\), define
\[
\begin{aligned}
Z_{aj}^{\mathrm B}
&=\sum_{i=1}^{h'}
\mathbf1\{X_i^{\mathrm o}=j,\ A_i^{\mathrm o}=a,\ Y_i^{\mathrm o}=1\},\\
K_{aj}^{\mathrm B}
&=\sum_{i=1}^{h''}
\mathbf1\{X_i^{\mathrm m}=j,\ A_i^{\mathrm m}=a\},\\
H_{aj}^{\mathrm B}
&=\frac{Z_{aj}^{\mathrm B}}{u_{\mathrm B}}
\left(1+\frac{K_{1-a,j}^{\mathrm B}}{K_{aj}^{\mathrm B}+1}\right),\\
F_{\mathrm B}(h',h'';\mathcal D)
&=\operatorname{clip}_{[-1,1]}
\left(\sum_{j=1}^d(H_{1j}^{\mathrm B}-H_{0j}^{\mathrm B})\right).
\end{aligned}
\]

\item Output the deterministic finite average
\[
T^{\mathrm B}(\mathcal D)
=
\sum_{h'=0}^{h_{\mathrm o}}\sum_{h''=0}^{h_{\mathrm m}}
\exp(-u_{\mathrm B}-t_{\mathrm B})
\frac{u_{\mathrm B}^{h'}}{h'!}
\frac{t_{\mathrm B}^{h''}}{h''!}
F_{\mathrm B}(h',h'';\mathcal D).
\]
The weights are the probabilities of independent Poisson prefix lengths with means \(u_{\mathrm B}\) and \(t_{\mathrm B}\); prefix pairs exceeding either pool length have output zero.
\end{enumerate}
\end{algorithmv}

\begin{theoremv}{thm:uniform-baseline}[Uniform inverse-count risk guarantee]
There exists a universal numerical constant \(C>0\), independent of \(n,m,d,\epsilon\), such that, for every experiment satisfying
\begin{itemize}
\item \textbf{(Sample sizes.)} \(n,m\) are nonnegative integers with \(n\ge 1\).
\item \textbf{(Alphabet size.)} \(d\) is an integer with \(d\ge 2\).
\item \textbf{(Overlap parameter.)} \(\epsilon\) is real and \(0<\epsilon\le 1/4\).
\end{itemize}
the baseline \(T^{\mathrm B}\) of \cref{obj:def:baseline} is measurable and satisfies
\[
-1\le T^{\mathrm B}(\mathcal D)\le 1
\]
for every data realization \(\mathcal D\). For every observable law \(P\in\mathcal M_{d,\epsilon}\), its squared-error risk satisfies
\[
\mathbb E_P\!\left[
  \bigl(T^{\mathrm B}(\mathcal D)-\tau(P)\bigr)^2
\right]
\le
C\min\!\left\{
1,\,
\frac{1}{n\epsilon}
+
\left(\frac{d}{(n+m)\epsilon}\right)^2
\right\}.
\]
The expectation is over the two-channel data under \(P\); the rule ignores the independent uniform randomization seed.
\end{theoremv}

\begin{definitionv}{synth_8}[Allowed experiment sequences]
An experiment sequence \(\mathbf v\) is an arbitrary sequence of public experiment indices, indexed by the nonnegative integers:
\[
\mathbf v=\bigl((n_k,m_k,d_k,\epsilon_k)\bigr)_{k=0}^{\infty}.
\]
For every \(k\), the complete-record size \(n_k\), auxiliary size \(m_k\), and alphabet size \(d_k\) are integers satisfying
\[
n_k\ge 1,\qquad m_k\ge 0,\qquad d_k\ge 2,
\]
and the public overlap floor \(\epsilon_k\) is a real number satisfying
\[
0<\epsilon_k\le \frac14.
\]
\end{definitionv}

\begin{definitionv}{def:phases}[Resource phases \(\mathcal C,\mathcal O,\mathcal I\)]
Define
\[
\begin{aligned}
\mathcal C
&=\left\{\mathbf v:
R(n_k,m_k,d_k,\epsilon_k)\longrightarrow 0\right\},\\
\mathcal O
&=\left\{\mathbf v:
n_k\epsilon_k\longrightarrow\infty,\quad
\limsup_{k\to\infty}
\frac{R(n_k,m_k,d_k,\epsilon_k)}{b(n_k,\epsilon_k)}
<\infty\right\},\\
\mathcal I
&=\left\{\mathbf v:
\frac{R(n_k,m_k,d_k,\epsilon_k)}
{R(n_k,0,d_k,\epsilon_k)}
\longrightarrow 0\right\}.
\end{aligned}
\]
Each set ranges over all allowed sequences of public experiment indices
\[
\mathbf v=((n_k,m_k,d_k,\epsilon_k))_{k\ge1}.
\]
For each such sequence, the rare-arm information scale and total marginal-information sample size are
\[
S_k=n_k\epsilon_k,\qquad N_k=n_k+m_k.
\]
Write \(\ell_k\) for the regularized logarithm \(\ell\) evaluated at the \(k\)-th experiment.
\end{definitionv}

\begin{theoremv}{thm:annotation-resource-phases}[Exact annotation resource phases]
Let \(\mathbf v=((n_k,m_k,d_k,\epsilon_k))_{k\ge0}\) be any sequence of allowed public experiment indices, and put
\[
S_k=n_k\epsilon_k,\qquad
N_k=n_k+m_k,\qquad
\ell_k=\log(\exp(1)+S_k).
\]
For the phases defined in \cref{obj:def:phases}, the following equivalences hold as \(k\to\infty\):
\[
\begin{aligned}
\mathbf v\in\mathcal C
&\Longleftrightarrow
S_k\longrightarrow\infty
\quad\text{and}\quad
d_k=o(N_k\epsilon_k\ell_k),\\
\mathbf v\in\mathcal O
&\Longleftrightarrow
S_k\longrightarrow\infty
\quad\text{and}\quad
d_k=O\!\left(\frac{N_k\epsilon_k\ell_k}{\sqrt{S_k}}\right),\\
\mathbf v\in\mathcal I
&\Longleftrightarrow
S_k\longrightarrow\infty,\quad
\frac{d_k}{\sqrt{S_k}\ell_k}\longrightarrow\infty,\quad
\frac{N_k/n_k}{\max\{1,d_k/(S_k\ell_k)\}}\longrightarrow\infty.
\end{aligned}
\]
If \(m_k=0\) for every \(k\), then
\[
\begin{aligned}
\mathbf v\in\mathcal C
&\Longleftrightarrow
S_k\longrightarrow\infty
\quad\text{and}\quad d_k=o(S_k\ell_k),\\
\mathbf v\in\mathcal O
&\Longleftrightarrow
S_k\longrightarrow\infty
\quad\text{and}\quad d_k=O(\sqrt{S_k}\ell_k),
\end{aligned}
\qquad
\mathbf v\notin\mathcal I.
\]
Whenever \(S_k\to\infty\), oracle membership also has the equivalent budget characterization
\[
\mathbf v\in\mathcal O
\quad\Longleftrightarrow\quad
\frac{d_k\sqrt{S_k}}{\epsilon_k\ell_k}=O(N_k).
\]

Write \(R\) for the minimax risk in \cref{obj:def:risk}. For every \(\mathbf v\in\mathcal O\), there exist constants \(c,C>0\) such that, for all sufficiently large \(k\),
\[
\frac{c}{S_k}
\le R(n_k,m_k,d_k,\epsilon_k)
\le \frac{C}{S_k}.
\]
Moreover, there exists \(c>0\) such that, for all sufficiently large \(k\) and every \(m'\in\mathbb N\),
\[
\frac{c}{S_k}\le R(n_k,m',d_k,\epsilon_k).
\]

For a single experiment, let \(n,m,d\in\mathbb N\) and \(\epsilon\in\mathbb R\) satisfy
\begin{itemize}
\item \textbf{(Complete-record budget.)} \(n\ge1\).
\item \textbf{(Alphabet size.)} \(d\ge2\).
\item \textbf{(Overlap range.)} \(0<\epsilon\le1/4\).
\end{itemize}
Put
\[
N=n+m,\qquad S=n\epsilon,\qquad
\ell=\log(\exp(1)+S),\qquad
b(n,\epsilon)=\min\{1,S^{-1}\}.
\]
There exists a universal constant \(c>0\) such that, for every such experiment,
\[
N\epsilon\ell\le d
\quad\Longrightarrow\quad
c\le R(n,m,d,\epsilon)\le1.
\]
There also exist universal constants \(c,C>0\) such that, for every such experiment,
\[
\frac{d\sqrt{n\epsilon}}{\epsilon\ell}\le N
\quad\Longrightarrow\quad
c\,b(n,\epsilon)\le R(n,m,d,\epsilon)\le C\,b(n,\epsilon).
\]
\end{theoremv}

\begin{definitionv}{synth_7}[Bernoulli law]
For \(p\in[0,1]\), \(\operatorname{Bernoulli}(p)\) is the probability law on \(\{0,1\}\) assigning probability \(p\) to \(1\) and probability \(1-p\) to \(0\).
\end{definitionv}

\begin{lemmav}{lem:causal-realization}[Causal realization and identification]
Let \(d\) be a natural number, \(\epsilon\) a real number, and \(P\) a probability law of the complete observation \((X,A,Y)\). Suppose:
\begin{itemize}
\item \textbf{(Covariate alphabet.)} \(d\ge 2\).
\item \textbf{(Overlap range.)} \(0<\epsilon\le 1/4\).
\item \textbf{(Model membership.)} \(P\in\mathcal M_{d,\epsilon}\), as defined in \cref{obj:def:model}.
\end{itemize}
Then \(P\) admits a joint probability law \(H\) of \((X,A,Y,Y_0,Y_1)\), with binary potential outcomes \(Y_0,Y_1\), whose \((X,A,Y)\) marginal equals \(P\) and which satisfies \cref{obj:ass:consistency,obj:ass:exchangeability}. Call precisely these extensions compatible.

One compatible extension is given by
\[
\Pr_H(X=j)=p_j(P),
\]
and, for every occupied cell \(j\), by the conditional law
\[
\mathcal L_H(A,Y_0,Y_1\mid X=j)
=
\operatorname{Bernoulli}\!\bigl(e_j(P)\bigr)
\otimes
\operatorname{Bernoulli}\!\bigl(\mu_{0j}(P)\bigr)
\otimes
\operatorname{Bernoulli}\!\bigl(\mu_{1j}(P)\bigr),
\qquad Y=Y_A.
\]
For every compatible extension \(H\), the causal contrast equals the target defined in \cref{obj:def:target}:
\[
\mathbb E_H[Y_1-Y_0]=\tau(P).
\]
\end{lemmav}

\begin{lemmav}{lem:poisson-inverse-moments}[Inverse Poisson count moments]
Let \(\lambda\ge 0\), let \(K\sim\operatorname{Poi}(\lambda)\) have the Poisson law with mean \(\lambda\), and define
\[
g(K)=\frac{1}{K+1},
\qquad C_0=2^{16}.
\]
Then
\[
\mathbb E[g(K)]
=
\begin{cases}
1,&\lambda=0,\\
\dfrac{1-e^{-\lambda}}{\lambda},&\lambda>0,
\end{cases}
\qquad
\mathbb E[g(K)^2]\le \frac{C_0}{(1+\lambda)^2},
\qquad
\operatorname{Var}(g(K))\le \frac{C_0\lambda}{(1+\lambda)^4}.
\]
\end{lemmav}

\begin{lemmav}{lem:inverse-count-arm-risk}[Inverse-count arm risk bounds]
There exists a universal numerical constant \(C>0\) such that the following holds.
\begin{itemize}
\item \textbf{(Parameters.)} Let \(d\ge 2\) be an integer, \(0<\epsilon\le 1/4\), \(a\in\{0,1\}\), and \(u,t>0\).
\item \textbf{(Population.)} Let \(P\in\mathcal M_{d,\epsilon}\), as defined in \cref{obj:def:model}. For \(j=1,\ldots,d\) and \(b\in\{0,1\}\), write
\[
p_j(P)=P(X=j),\qquad
s_{bj}(P)=P(X=j,A=b),\qquad
q_{bj}(P)=P(X=j,A=b,Y=1),
\]
and define the outcome-success conditional mean by
\[
\mu_{bj}(P)=
\begin{cases}
q_{bj}(P)/s_{bj}(P),&s_{bj}(P)>0,\\
0,&s_{bj}(P)=0.
\end{cases}
\]
\item \textbf{(Poisson counts.)} On a common probability space, let \(Z_j,K_j,W_j\) be measurable counts satisfying
\[
Z_j\sim\operatorname{Poi}\!\bigl(uq_{aj}(P)\bigr),\qquad
K_j\sim\operatorname{Poi}\!\bigl(ts_{aj}(P)\bigr),\qquad
W_j\sim\operatorname{Poi}\!\bigl(ts_{1-a,j}(P)\bigr),
\]
where \(\operatorname{Poi}(\lambda)\) denotes the Poisson law with mean \(\lambda\). The entire collection of \(3d\) counts is mutually independent.
\end{itemize}
Define the cell contributions and the covariate-weighted arm mean by
\[
H_j=\frac{Z_j}{u}\left(1+\frac{W_j}{K_j+1}\right),
\qquad
\psi_a=\sum_{j=1}^{d}p_j(P)\mu_{aj}(P).
\]
Then
\[
\left|\mathbb E\!\left[\sum_{j=1}^{d}H_j\right]-\psi_a\right|
\le
\min\left\{1,\frac{d}{\exp(1)t\epsilon}\right\},
\qquad
\operatorname{Var}\!\left(\sum_{j=1}^{d}H_j\right)
\le
C\left(\frac{1}{u\epsilon}+\frac{1}{t\epsilon}\right).
\]
\end{lemmav}

\begin{lemmav}{lem:finite-prefix-transfer}[Finite prefix risk transfer]
Let \(I\) be a finite index set. Suppose the following conditions hold.
\begin{itemize}
\item \textbf{(Pools.)} For each \(i\in I\), let \(P_i\) be a probability law on an arbitrary measurable observation space, and let \(h_i\geq 1\) be an integer pool length. Observations within pool \(i\) are iid with law \(P_i\), and the pools are mutually independent.
\item \textbf{(Statistic.)} Let \(T\) be a measurable real-valued statistic on families of finite ordered prefixes, including empty prefixes, such that
\[
-1\leq T(s)\leq 1
\]
for every prefix family \(s\).
\item \textbf{(Target and risk.)} Let \(\vartheta\in[-1,1]\) be the target and \(A\in\mathbb R\) the risk bound. Let \(J_i\sim\operatorname{Poi}(h_i/8)\) be mutually independent prefix lengths, independent of all observations. For independent infinite iid observation sequences with respective laws \(P_i\), suppose
\[
\mathbb E\!\left[
\left(T\bigl((s_{i,1:J_i})_{i\in I}\bigr)-\vartheta\right)^2
\right]\leq A.
\]
\end{itemize}
For finite ordered pools \(s=(s_{i,1:h_i})_{i\in I}\), define the prefix-averaged statistic by
\[
T_{\mathrm{fin}}(s)
=
\sum_{\substack{(j_i)_{i\in I}\\0\leq j_i\leq h_i\ \forall i}}
\left(
\prod_{i\in I}
e^{-h_i/8}\frac{(h_i/8)^{j_i}}{j_i!}
\right)
T\bigl((s_{i,1:j_i})_{i\in I}\bigr).
\]
This is the conditional average over the prefix lengths with output zero whenever some \(J_i>h_i\). It defines a total deterministic measurable rule, takes values in \([-1,1]\) for every finite pool family, and satisfies
\[
\mathbb E\!\left[
\left(T_{\mathrm{fin}}\bigl((s_{i,1:h_i})_{i\in I}\bigr)
-\vartheta\right)^2
\right]
\leq A+\sum_{i\in I}e^{-h_i},
\]
where the expectation is under the independent finite iid pools. For observation spaces equipped with their Borel sigma-algebras, the rule is Borel measurable.
\end{lemmav}

\begin{lemmav}{lem:chebyshev-factorial-certificate}[Chebyshev factorial moment certificate]
Let \(L\) be an integer with \(L\ge2\). Let \(T_L\) be the Chebyshev polynomial specified by
\[
T_0(x)=1,\qquad T_1(x)=x,\qquad
T_{h+1}(x)=2xT_h(x)-T_{h-1}(x)\quad(h\ge1),
\]
and define
\[
C_L(z)=T_L(1-2z).
\]
Define the polynomials \(E_L\) and \(G_L\) by the identities
\[
E_L(z)=\frac{1-C_L(z)}{2L^2z},
\qquad
G_L(z)=\frac{1-E_L(z)}{z}
\quad(z\ne0),
\]
with their removable values at zero. Then
\[
0\le E_L(z)\le\min\{1,(L^2z)^{-1}\}
\quad(0<z\le1),
\qquad E_L(0)=1.
\]
Moreover,
\[
\deg G_L\le L-2,\qquad
G_L(z)=\sum_{h=0}^{L-2}g_hz^h,\qquad
\sum_{h=0}^{L-2}|g_h|\le14^L,
\]
where \(g_h\) is the coefficient of \(z^h\) in \(G_L\).

For real parameters \(t,B,z\), define
\[
D=tB,
\]
and suppose that
\begin{itemize}
\item \textbf{(Positive scales.)} \(t>0\) and \(B>0\).
\item \textbf{(Nonnegative argument.)} \(z\ge0\).
\item \textbf{(Degree calibration.)} \(D\ge L\).
\end{itemize}
Let \(K\sim\operatorname{Poi}(tzB)\), the Poisson law with mean \(tzB\), and define
\[
(K)_0=1,\qquad
(K)_h=\prod_{i=0}^{h-1}(K-i)\quad(h\ge1),
\qquad
\widehat G=\sum_{h=0}^{L-2}g_h\frac{(K)_h}{D^h}.
\]
Then
\[
\mathbb E[\widehat G]=G_L(z),
\qquad
\mathbb E[\widehat G^2]\le392^L(1+z)^{2L}.
\]
\end{lemmav}

\begin{lemmav}{lem:uniform-hybrid-arm-risk}[Uniform hybrid arm risk]
There exists a numerical constant \(C>0\) such that the following holds for every \(n,m,d\in\mathbb N\), overlap parameter \(\epsilon\in\mathbb R\), observable law \(P\), arm \(a\in\{0,1\}\), and pool intensities \(u,t_p,t\in\mathbb R\).
\begin{itemize}
\item \textbf{(Experiment.)} The indices satisfy
\[
n\ge1,\qquad d\ge2,\qquad 0<\epsilon\le\frac14,\qquad
S=n\epsilon\ge\exp(4096),\qquad N=n+m.
\]
\item \textbf{(Model.)} The law satisfies \(P\in\mathcal M_{d,\epsilon}\), as defined in \cref{obj:def:model}.
\item \textbf{(Intensities.)} The pool intensities satisfy
\[
u\ge\frac n{32},\qquad t_p\ge\frac N{64},\qquad
t\ge\frac N{64},\qquad \frac13\le\frac t{t_p}\le3.
\]
\item \textbf{(Calibration.)} Set
\[
L=\left\lfloor\frac{\log S}{1024}\right\rfloor,\qquad
H_0=2^{20},\qquad B=\frac{H_0L}{\min\{t_p,t\}},\qquad
k_0=\left\lfloor\frac{t_pB}{4}\right\rfloor.
\]
\item \textbf{(Poisson pools.)} On any probability space, let
\(\{Z_{bj},J_{bj},K_{bj}:b\in\{0,1\},\,1\le j\le d\}\)
be jointly independent measurable counts with
\[
Z_{bj}\sim\operatorname{Poi}\!\left(uq_{bj}(P)\right),\qquad
J_{bj}\sim\operatorname{Poi}\!\left(t_ps_{bj}(P)\right),\qquad
K_{bj}\sim\operatorname{Poi}\!\left(ts_{bj}(P)\right).
\]
Here \(\operatorname{Poi}(\lambda)\) is the Poisson law with mean \(\lambda\), and the population quantities are
\[
q_{bj}(P)=P(X=j,A=b,Y=1),\qquad
s_{bj}(P)=\sum_{y\in\{0,1\}}P(X=j,A=b,Y=y),
\]
\[
p_j(P)=\sum_{b\in\{0,1\}}s_{bj}(P),\qquad
\mu_{bj}(P)=
\begin{cases}
q_{bj}(P)/s_{bj}(P),&s_{bj}(P)>0,\\
0,&s_{bj}(P)=0.
\end{cases}
\]
\end{itemize}
Using the monomial coefficients \(g_h\) of the reciprocal-approximation polynomial \(G_L\) in the estimator construction, form
\[
\widehat G_{aj}
=\sum_{h=0}^{L-2}g_h\frac{(K_{aj})_h}{(tB)^h},
\qquad
(K)_0=1,\qquad
(K)_h=\prod_{r=0}^{h-1}(K-r)\quad(h\ge1),
\]
and
\[
V_{aj}=
\begin{cases}
\displaystyle
\frac{Z_{aj}}u
\left(1+\frac{K_{1-a,j}\widehat G_{aj}}{tB}\right),
&J_{aj}\le k_0,\\[6pt]
\displaystyle
\frac{Z_{aj}}u
\left(1+\frac{K_{1-a,j}}{K_{aj}+1}\right),
&J_{aj}>k_0.
\end{cases}
\]
Then
\[
\mathbb E\!\left[
\left(\sum_{j=1}^dV_{aj}
-\sum_{j=1}^dp_j(P)\mu_{aj}(P)\right)^2
\right]
\le
C\left\{\frac1S+
\left(\frac d{N\epsilon L}\right)^2\right\}.
\]
The same constant \(C\) applies to every choice of indices, law, arm, intensities, and probability space satisfying these conditions.
\end{lemmav}

\begin{lemmav}{lem:pinsker-finite}[Pinsker inequality on finite alphabets]
Let \(\nu,\nu'\) be probability measures on a finite set. Define their total variation distance and natural-logarithm Kullback--Leibler divergence by
\[
\operatorname{TV}(\nu,\nu')
=\frac12\sum_z|\nu(z)-\nu'(z)|,
\qquad
D(\nu\Vert\nu')
=
\begin{cases}
\displaystyle\sum_z \nu(z)\log\frac{\nu(z)}{\nu'(z)},
& \nu'(z)=0\Longrightarrow \nu(z)=0\ \text{for every }z,\\
+\infty,& \text{otherwise}.
\end{cases}
\]
Each summand with \(\nu(z)=0\) is defined to be zero. Then, with the inequality interpreted in the extended nonnegative reals,
\[
\operatorname{TV}(\nu,\nu')
\le \sqrt{\frac{D(\nu\Vert\nu')}{2}}.
\]
\end{lemmav}

\begin{lemmav}{lem:markov-polynomial}[Markov inequality for polynomials]
Let \(L\) be a nonnegative integer, \(g\) a real polynomial of degree at most \(L\), and \(M_b\in\mathbb R\). Suppose that
\[
|g(x)|\le M_b \qquad \text{for every } x\in[-1,1].
\]
Then
\[
|g'(x)|\le L^2 M_b \qquad \text{for every } x\in[-1,1].
\]
\end{lemmav}

\begin{lemmav}{lem:finite-polynomial-duality}[Finite polynomial approximation duality]
Let \(v_-<v_+\), let \(g:[v_-,v_+]\to\mathbb R\) be continuous, and let \(L\ge0\) be an integer. There exist nodes
\[
v_-\le v_1<\cdots<v_{L+2}\le v_+
\]
and real weights \(w_1,\ldots,w_{L+2}\) such that
\[
\sum_{i=1}^{L+2}|w_i|=1,
\qquad
\sum_{i=1}^{L+2}w_i v_i^h=0
\quad\text{for every integer }0\le h\le L,
\]
and
\[
\sum_{i=1}^{L+2}w_i g(v_i)
=
\inf_{Q:\,\deg Q\le L}
\sup_{v\in[v_-,v_+]}|g(v)-Q(v)|,
\]
where the infimum ranges over real polynomials of degree at most \(L\). Equivalently, the finite signed measure
\[
\sigma=\sum_{i=1}^{L+2}w_i\delta_{v_i}
\]
has support on at most \(L+2\) points and satisfies
\[
\|\sigma\|_{\mathrm{TV}}=1,
\qquad
\int v^h\,d\sigma(v)=0
\quad\text{for every integer }0\le h\le L,
\qquad
\int g\,d\sigma
=
\inf_{Q:\,\deg Q\le L}
\sup_{v\in[v_-,v_+]}|g(v)-Q(v)|.
\]
Here \(\delta_{v_i}\) is the unit point mass at \(v_i\), and \(\|\sigma\|_{\mathrm{TV}}=\sum_{i=1}^{L+2}|w_i|\) is the total mass of the variation measure of \(\sigma\).
\end{lemmav}

\begin{lemmav}{lem:shrinking-cone-dual}[Finite reciprocal-gap certificate]
Let the parameters satisfy:
\begin{itemize}
\item \textbf{(Degree.)} \(L\) is an integer with \(L\ge 2\).
\item \textbf{(Interval.)} \(B>0\).
\item \textbf{(Overlap.)} \(0<\epsilon\le 1/4\).
\end{itemize}
Set
\[
\alpha=\frac{B}{100L^2}.
\]
There exist strictly increasing nodes \(0\le v_0<\cdots<v_{L+1}\le B\) and real weights \(c_0,\ldots,c_{L+1}\) defining the finite signed measure
\[
\sigma=\sum_{i=0}^{L+1}c_i\delta_{v_i},
\]
where \(\delta_{v_i}\) is the unit point mass at \(v_i\), such that
\[
\|\sigma\|_{\mathrm{TV}}
=\sum_{i=0}^{L+1}|c_i|=1,
\qquad
\int v^h\,d\sigma(v)=0
\quad\text{for every integer }0\le h\le L,
\]
and
\[
\int\frac{\alpha}{\alpha+(1-\epsilon)v}\,d\sigma(v)
\ge\frac18.
\]
Here \(\|\sigma\|_{\mathrm{TV}}\) denotes the total mass of the variation measure of \(\sigma\).
\end{lemmav}

\begin{definitionv}{synth_9}[Rare-label probability tables]
For \(d\ge2\), \(0<\epsilon\le1/4\), and \(0<\delta\le1/8\), define \(P_1\) and \(P_0\) on \(\{1,\ldots,d\}\times\{0,1\}^2\) by
\[
\begin{aligned}
P_1(X=j,A=1,Y=1)&=\frac{\epsilon}{2}\left(\frac12+\delta\right),&
P_1(X=j,A=1,Y=0)&=\frac{\epsilon}{2}\left(\frac12-\delta\right),\\
P_0(X=j,A=1,Y=1)&=\frac{\epsilon}{2}\left(\frac12-\delta\right),&
P_0(X=j,A=1,Y=0)&=\frac{\epsilon}{2}\left(\frac12+\delta\right),
\end{aligned}
\qquad j\in\{1,2\}.
\]
For both tables, set
\[
P_1(X=j,A=0,Y=y)=P_0(X=j,A=0,Y=y)=\frac{1-\epsilon}{4},
\qquad j\in\{1,2\},\quad y\in\{0,1\},
\]
and assign zero probability to every record with \(j>2\).
Thus \(P_1\) is the positive-sign law and \(P_0\) the negative-sign law of \(\cref{obj:thm:rare-label-floor}\). In the rare-label branch of \(\cref{obj:def:prior-handle}\), use these tables with
\[
\delta=\frac{1}{16\sqrt{\max\{1,n\epsilon\}}},
\qquad
\Pi_1=\delta_{P_1},\qquad \Pi_0=\delta_{P_0},
\]
where \(\delta_P\) denotes the Dirac probability measure at the table \(P\).
\end{definitionv}

\begin{definitionv}{synth_6}[Total-variation measure]
For the selected finite signed measure \(\sigma\), let \(|\sigma|\) denote its total-variation measure. For every measurable set \(E\),
\[
|\sigma|(E)=\sup_{\{E_1,\ldots,E_r\}}\sum_{i=1}^{r}|\sigma(E_i)|,
\]
where the supremum ranges over all finite measurable partitions of \(E\). In particular, if \(\sigma=\sum_{i=1}^{r}c_i\delta_{v_i}\) with distinct atoms \(v_i\), then
\[
|\sigma|=\sum_{i=1}^{r}|c_i|\delta_{v_i}.
\]
\end{definitionv}

\begin{algorithmv}{def:prior-handle}[Two testing prior construction]
The testing priors \(\Pi_1,\Pi_0\) are measures on complete-record probability tables, defined for public parameters \(n,m,d\in\mathbb N\) and \(\epsilon\in\mathbb R\) satisfying
\begin{itemize}
\item \textbf{(Complete-record budget.)} \(n\ge1\).
\item \textbf{(Alphabet size.)} \(d\ge2\).
\item \textbf{(Overlap range.)} \(0<\epsilon\le1/4\).
\end{itemize}
Their construction is as follows.
\begin{enumerate}
\item Put
\[
S=n\epsilon,\qquad
\ell=\log(\exp(1)+S),\qquad
N=n+m,\qquad
x=\frac d{N\epsilon\ell}.
\]
If \(S<\exp(4096)\) or \(x^2\le S^{-1}\), set
\[
\delta=\frac1{16\sqrt{\max\{1,S\}}}.
\]
For \(r\in\{0,1\}\), define the two-cell table by
\[
P_r(X=j,A=a,Y=y)=
\begin{cases}
\displaystyle
\frac12\epsilon^a(1-\epsilon)^{1-a}
\left(\frac12+(2r-1)a\delta\right)^y
\left(\frac12-(2r-1)a\delta\right)^{1-y},
&j\in\{1,2\},\\[4pt]
0,&3\le j\le d,
\end{cases}
\qquad a,y\in\{0,1\}.
\]
Output
\[
\Pi_1=\delta_{P_1},\qquad \Pi_0=\delta_{P_0}.
\]
Otherwise continue with the following steps.

\item Set
\[
\begin{aligned}
u&=128n,\qquad w=128N,\qquad M=u+w,\\
L&=\lceil8\ell\rceil,\qquad
B=\frac L{1000M},\qquad
\alpha=\frac B{100L^2},\\
b_0&=\frac\alpha\epsilon,\qquad
K_*=\min\left\{d-1,\left\lfloor\frac1{100b_0}\right\rfloor\right\}.
\end{aligned}
\]

\item Fix a witness supplied by \cref{obj:lem:shrinking-cone-dual}, represented as
\[
\sigma=\sum_{i=1}^{L+2}\omega_i\delta_{v_i},
\]
with defining properties
\[
\begin{gathered}
0\le v_1<\cdots<v_{L+2}\le B,\qquad
\sum_{i=1}^{L+2}|\omega_i|=1,\\
\sum_{i=1}^{L+2}\omega_i v_i^h=0
\quad(h=0,\ldots,L),\\
\sum_{i=1}^{L+2}\omega_i
\frac{\alpha}{\alpha+(1-\epsilon)v_i}\ge\frac18.
\end{gathered}
\]
At each atom of positive variation mass, define the polar sign
\[
h(v)=\frac{\sigma(\{v\})}{|\sigma|(\{v\})}\in\{-1,1\}.
\]

\item Define the raw covariate and arm masses by
\[
w_*(v)=b_0+v,\qquad
S_1(v)=\alpha+(1-\epsilon)v,\qquad
S_0(v)=(1-\epsilon)b_0+\epsilon v,
\qquad 0\le v\le B.
\]
For each rare cell \(i=1,\ldots,K_*\), independently draw a latent coordinate \(V_i\) and a branch flag \(c_i\in\{0,1\}\), with \(c_i=1\) denoting the core branch, according to
\[
\begin{aligned}
\Pr(V_i\in dv,\ c_i=1)
&=\frac{\alpha}{S_1(v)}|\sigma|(dv),\\
\Pr(V_i=0,\ c_i=0)
&=1-\int\frac{\alpha}{S_1(v)}|\sigma|(dv).
\end{aligned}
\]
The filler and core branches remain distinct when the core measure has an atom at zero. For \(r\in\{0,1\}\), set
\[
\mu_1^{(r)}(v,c)=
\begin{cases}
\bigl(1+(2r-1)h(v)\bigr)/2,&c=1,\\
0,&c=0,
\end{cases}
\qquad
\mu_0^{(r)}(v,c)=0.
\]
Both priors use raw arm masses \(S_1(v),S_0(v)\).

\item Let \(V\) have the common latent-coordinate marginal just defined, and put
\[
\bar w=\mathbb E[w_*(V)],\qquad
p_*=1-K_*\bar w,\qquad
Q=p_*+\sum_{i=1}^{K_*}w_*(V_i).
\]
Append a reservoir cell with raw covariate mass \(p_*\), propensity \(1/2\), and both outcome means zero. Set the remaining cells null.

\item For \(r\in\{0,1\}\), define the normalized complete-record table by
\[
P_r(X=j,A=a,Y=y)=
\begin{cases}
\displaystyle
\frac{S_a(V_j)}Q
\left[
y\mu_a^{(r)}(V_j,c_j)
+(1-y)\bigl(1-\mu_a^{(r)}(V_j,c_j)\bigr)
\right],
&1\le j\le K_*,\\[6pt]
\displaystyle\frac{p_*}{2Q}\mathbf1\{y=0\},
&j=K_*+1,\\[6pt]
0,&K_*+1<j\le d,
\end{cases}
\qquad a,y\in\{0,1\}.
\]
Output the pushforward priors
\[
\Pi_1=\operatorname{Law}(P_1),\qquad
\Pi_0=\operatorname{Law}(P_0)
\]
of the common finite product latent distribution under the respective table maps.
\end{enumerate}
\end{algorithmv}

\begin{lemmav}{lem:normalized-affine-testing}[Normalized affine testing priors]
There exists a numerical constant \(c>0\) such that the following holds for every \(n,m,d\in\mathbb N\) and \(\epsilon\in\mathbb R\). Put
\[
S=n\epsilon,\qquad N=n+m,\qquad
\ell=\log(\exp(1)+S),\qquad x=\frac{d}{N\epsilon\ell},
\]
where \(S\) is the rare-arm information scale, \(N\) is the total marginal-information sample size, \(\ell\) is the regularized logarithmic scale, and \(x\) is the normalized many-cell difficulty. Suppose that:
\begin{itemize}
\item \textbf{(Public indices.)} \(n\ge1\) and \(d\ge2\).
\item \textbf{(Overlap parameter.)} \(0<\epsilon\le1/4\).
\item \textbf{(Information scale.)} \(S\ge\exp(4096)\).
\item \textbf{(Affine regime.)} \(x^2>S^{-1}\).
\end{itemize}
For every signed certificate \(\sigma\) satisfying the certificate conditions at the affine tuning parameters in \cref{obj:def:prior-handle}, and every latent configuration in that construction, the normalized laws under both hypotheses belong to \(\mathcal M_{d,\epsilon}\) of \cref{obj:def:model}, and their treatment--covariate marginal tables \(P_{XA}\) agree. For each such \(\sigma\), the two product priors obtained by pushing the common product latent distribution through the normalized-law construction also induce exactly the same distribution of \(P_{XA}\).

The selected priors \(\Pi_1\) and \(\Pi_0\) in \cref{obj:def:prior-handle} induce exactly the same distribution of \(P_{XA}\), and
\[
R(n,m,d,\epsilon)\ge c\min\{1,x^2\}.
\]
Here \(R(n,m,d,\epsilon)\) is the minimax squared-error risk in \cref{obj:def:risk}, with the infimum over all clipped Borel randomized rules in the original fixed-size two-channel experiment.
\end{lemmav}
